\documentclass[11pt]{article}
\usepackage{amsmath,amssymb,amsthm}
\usepackage{tikz-cd}

\theoremstyle{plain}
\newtheorem{theorem}{Theorem}[section]
\newtheorem{lemma}[theorem]{Lemma}
\theoremstyle{definition}
\newtheorem{definition}[theorem]{Definition}
\newtheorem{example}[theorem]{Example}
\newtheorem*{convention}{Convention}

\newcommand{\kk}{k}
\newcommand{\RR}{\mathbb{R}}
\newcommand{\Vec}{\mathbf{Vec}}
\newcommand{\Hloc}[2]{H^{#1}_{#2}}
\DeclareMathOperator{\Dgm}{Dgm}
\DeclareMathOperator{\link}{lk}

\title{A Sample Paper on Persistence}
\author{Ada Lovelace\thanks{Somewhere University} \and Emmy Noether}
\date{}

\begin{document}
\maketitle

\begin{abstract}
We study persistence modules over the poset $(\RR,\le)$ and prove a stability theorem (Theorem~\ref{thm:stab}).
\end{abstract}

\section{Preliminaries}\label{sec:prelim}

Let $\kk$ be a field and $S=\kk[x_1,\dots,x_n]$ the polynomial ring with its standard grading.
Let $\Delta$ be an abstract simplicial complex on the vertex set $[n]=\{1,\dots,n\}$.
For $W\subseteq[n]$ write $\Delta_W=\{\sigma\in\Delta:\sigma\subseteq W\}$ for the induced subcomplex.
The \emph{Stanley--Reisner ideal} is
\begin{equation}\label{eq:sr}
I_\Delta=\big(x_{i_1}\cdots x_{i_r}:\{i_1,\dots,i_r\}\notin\Delta\big)\subseteq S.
\end{equation}
The \emph{vertex prime} $\mathfrak p_i=(x_j:j\neq i)$ contains $I_\Delta$ by \eqref{eq:sr}.

\section{Persistence modules}\label{sec:modules}

\begin{definition}[Persistence module]\label{def:pm}
Let $(P,\le)$ be a poset. A \emph{$P$-indexed persistence module} is a functor $M\colon P\to\Vec$.
We write $M_t$ for $M(t)$ and $\varphi_M(s,t)\colon M_s\to M_t$ for the structure maps, following~\cite[Definition~2.1]{chazal2016structure}.
Such an $M$ is \emph{pointwise finite-dimensional} \textup{(}pfd\textup{)} if every $M_t$ is finite-dimensional.
\end{definition}

\begin{definition}[Interleaving distance]\label{def:interleave}
For $\delta\ge 0$ define
\begin{align}
d_I(M,N) &:= \inf\{\delta\ge0: M,N \text{ are } \delta\text{-interleaved}\},\label{eq:di}\\
\rho(M) &= \sup_t \dim M_t. \notag
\end{align}
\end{definition}

\begin{convention}
Throughout, $\Hloc{q}{\mathfrak m}(-)$ denotes local cohomology with support in $\mathfrak m$.
\end{convention}

\begin{example}
Let $\Delta$ be the path $1-2-3$. Then $\link_\Delta(2)=\{\varnothing,\{1\},\{3\}\}$.
\end{example}


\section{Stability}\label{sec:stab}

\begin{theorem}[Stability]\label{thm:stab}
Let $M$ and $N$ be pfd persistence modules. If $M$ and $N$ are $\varepsilon$-interleaved then
$d_B(\Dgm(M),\Dgm(N))\le\varepsilon$.
\end{theorem}

\begin{proof}
Fix a vertex $i$ and let $j\in\{1,\dots,n\}$ be an index. Since $M$ and $N$ are pfd, by Definition~\ref{def:pm} and \cite[Thm.~4.4]{bauer2015induced} the claim follows; see also \S\ref{sec:prelim}.
\end{proof}

\begin{lemma}\label{lem:cases}
The number $\nu(c)$ satisfies
\[
\nu(c)=
\begin{cases}
1, & c=0,\\[2pt]
0, & c\ne 0.
\end{cases}
\]
\end{lemma}

The following diagram commutes:
\[
\begin{tikzcd}
A \arrow[r] \arrow[d] & B \arrow[d] \\
C \arrow[r] & D
\end{tikzcd}
\]

\begin{figure}[h]
\centering
\begin{tikzpicture}
\draw (0,0) -- (1,0) -- (0.5,1) -- cycle;
\end{tikzpicture}
\caption{A triangle.}\label{fig:tri}
\end{figure}

\section*{Acknowledgements}
Thanks to everyone.\footnote{Really everyone.}

\bibliographystyle{plain}
\bibliography{refs}
\end{document}
